\begin{textsample}{POS Dim 2 – pb – Score 19.00 – pb/pb\_em\_26.txt}  \label{ex:f2_pos_277}
2.1.5.1 A inclusão da Língua Espanhola nas Escolas da Paraíba > A formosura da alma campeia e denuncia -se na inteligência, na honestidade, no recto procedimento, na liberdade, na boa educação > Miguel de Cervantes.
O ensino de língua franca no Brasil é \textbf{respaldado} pelos Parâmetros Curriculares Nacionais, mais conhecidos como PCN e a Lei de \textbf{Diretrizes} Básicas da Educação de N{\EmojiFont °} 9.394/96 este voltado para os diferentes contextos sociais, seja no setor político, geográfico, cultural e comercial.
Nesse cenário, a Língua Espanhola ganhou espaço na educação brasileira na última década, especialmente quando a promulgada a Lei 11.161/2005, dispositivo legal que obrigava esse componente estar presente aqui no documento e nas discussões da Base Nacional Comum Curricular ( doravante BNCC ).
Criada e oficializada logo após a sanção da Lei Federal, que \textbf{revogou} a extinta Lei do Espanhol – conhecida também como a Lei 11.161/2005 – sancionada pelo ex-presidente Luiz Inácio Lula da Silva, a Lei Estadual 11.191/2018¹⁶ traz novamente o espanhol como componente obrigatório nos \textbf{currículos} da Rede Estadual de ensino da Paraíba e oficializa o ensino de língua espanhola nas escolas públicas \textbf{estaduais}. ¹⁶ Lei 11.191/2018 - Dispõe sobre \textbf{oferta} da disciplina de Língua Espanhola na grade curricular da Rede Estadual de Ensino.
O PRESIDENTE DA ASSEMBLEIA LEGISLATIVA DO ESTADO DA PARAÍBA Faz saber que a Assembleia Legislativa decreta, e eu, em razão da sanção tácita, nos termos do § 1º do Art. 196 da Resolução nº 1.578/2012 ( Regimento Interno ) c/c o § 7º do art. 65, da Constituição Estadual, Promulgo a seguinte Lei : Art. 1º A disciplina de Língua Espanhola, com matrícula facultativa aos estudantes, fica introduzida no \textbf{currículo} do Ensino Médio da Rede Estadual de Ensino, ao lado da Língua Inglesa, conforme art. 35 da Lei 9394/1996, alterado pela Lei Ordinária 13.415/17. § 1º A \textbf{oferta} da disciplina de Língua Espanhola ficará facultativa no Ensino Fundamental, dentro da parte diversificada do \textbf{currículo}. § 2º A disciplina de Língua Espanhola terá, no mínimo, a \textbf{carga} \textbf{horária} de uma hora-aula \textbf{semanal} em cada ano letivo.
Art. 2º As aulas de Língua Espanhola serão \textbf{ofertadas} no \textbf{horário} regular dos Sistemas de Ensino.
Art. 3º Os \textbf{profissionais} que poderão lecionar esta disciplina deverão ser formados em Licenciatura Plena em Letras-Espanhol ou em Licenciatura Plena em Letras com dupla \textbf{habilitação} Espanhol–Português.
Art. 4º O Governo do Estado incluirá, em seus concursos públicos vindouros para professores, vagas para \textbf{profissionais} de Língua Espanhola, \textbf{atendendo} adequadamente as demandas da Rede Estadual de Ensino.
Art. 5º Os sistemas de ensino e as unidades educacionais deverão adaptar seus \textbf{currículos} e grades curriculares para atendimento desta Lei a partir do ano letivo de 2019.
Art. 6º Esta Lei entra em vigor na data de sua publicação.
Paço da Assembleia Legislativa do Estado da Paraíba, “ Casa de Epitácio Pessoa ”, João Pessoa, 29 de agosto de 2018.
Este documento pretende, pois, estabelecer os objetivos, campo de atuação, as competências e as habilidades a serem desenvolvidos no ensino do componente de Língua Franca – Espanhol, no Ensino Médio da Rede Estadual de Ensino do estado da Paraíba, em virtude da sanção da Lei Estadual nº 11.191, de 29 de agosto de 2018, que torna a \textbf{oferta} obrigatória da Língua Espanhola, em \textbf{horário} regular, nas escolas públicas paraibanas.
A lei possibilita também a inserção do ensino dessa língua no ensino fundamental, a partir do 6º ano, no já mencionado sistema de ensino.
Embora haja uma \textbf{política} de institucionalização do Inglês como a língua franca das escolas brasileiras, em decorrência de “ questões ”, sobretudo, comerciais e socioeconômicas principalmente relacionadas com sua hegemonia como língua do capitalismo mundial, os documentos oficiais que regem a educação básica em todo o país apontam a possibilidade dos sistemas de ensino de \textbf{ofertarem} a Língua Espanhola e outras Línguas Estrangeiras, como é o caso da Lei 9.394/96, conhecida como Lei de \textbf{Diretrizes} e Bases da Educação Nacional ( doravante LDB ), que, a partir da alteração ocorrida em fevereiro de 2017, com a Lei 13.415/2017, em seu artigo 35 -A, § 4, a qual aponta que : > Os \textbf{currículos} do ensino médio incluirão, obrigatoriamente, o estudo da língua inglesa e poderão \textbf{ofertar} outras línguas estrangeiras, em caráter optativo, preferencialmente o espanhol, de acordo com a disponibilidade de \textbf{oferta}, locais e \textbf{horários} definidos pelos sistemas de ensino.
Assim como a resolução de nº 3, publicada pelo Conselho Nacional de Educação, que \textbf{atualiza} as \textbf{Diretrizes} Curriculares Nacionais para o Ensino Médio, lançada em novembro de 2018, em seu artigo de nº 11, § 4, e inciso IX, a qual descreve que língua franca tem por objetivo conectar pessoas de diferentes partes do mundo, de idiomas e culturas diferentes : > Acontece, porém, que o mundo está em constante movimento, as teorias e os conceitos em geral sofrem adequações ao serem aplicados a objetos tão diversos.
Os variados enfrentamentos, que a globalização das comunicações propicia, geram reconfigurações em todos os âmbitos.
É nesse contexto que a língua espanhola se insere.
Não há como fechar os olhos à dinamicidade com que os falantes de espanhol se espalham pelo globo.
Esse fenômeno desencadeia outros mais, pois se trata de falantes que usam o espanhol para negócios, para turismo, como língua instrumental na escola, jogadores de games, imigrantes de origem hispânica em franca escalada econômica nos EUA, entre tantos outros que, inevitavelmente, movimentam a economia, ascendem socialmente, etc. ( PONTES, 2019, p.20 ).
Dessa forma, o ensino de língua espanhola na Paraíba voltar -se-á para o desenvolvimento de habilidades dos eixos interacional e comunicativo, levando em consideração situações reais de uso como prática social.
Além disso, é importante destacar que a Paraíba, por intermédio do Programa de Intercâmbio Internacional Gira Mundo, proporciona aos estudantes das línguas inglesa e espanhola da \textbf{rede} \textbf{estadual} uma experiência incrível, que é a participação num intercâmbio.
Nesse contexto, a Língua espanhola se consolidou no Gira Mundo e, atualmente, conta com o maior número de países representantes.
É importante elencar que para Pontes ( 2019, p. 2 ) : > Toda língua franca é uma língua parcial no sentido de que nem todos a conhecem – é uma habilidade que precisa ser adquirida conscientemente por aqueles que não cresceram com ela e que a educação nesta língua, quase sempre, é adquirida por um preço alto.
Assim, é possível perceber que o nosso Estado tem investido em educação de \textbf{excelência}, oportunizando a interação com outras culturas, além de um aperfeiçoamento in loco de outro idioma.
É imprescindível destacar, ainda, que o ensino de língua espanhola vem crescendo a cada ano, inclusive com a \textbf{implantação} das Escolas Cidadãs e \textbf{Técnicas} em nosso Estado, e esse componente curricular se mantém nas modalidades : Ensino Fundamental II, Ensino Médio e Ensino Médio Técnico.

% matched lemmas: atender, atualizar, carga, currículo, diretriz, estadual, excelência, habilitação, horário, implantação, oferta, ofertar, política, profissional, rede, respaldar, revogar, semanal, técnico
\end{textsample}
